% Part: propositional-logic
% Chapter: syntax-and-semantics
% Section: formation-sequences

\documentclass[../../../include/open-logic-section]{subfiles}

\begin{document}

\olfileid[tr]{pl}{syn}{fseq}

\olsection{Oluşum Dizileri}

!!{formula}s için tümevarımlı bir tanım kullanmak ve !!{formula}s hakkında
tümevarımla özellikler kanıtlamak zarif ve verimli bir yaklaşımdır. Bununla
birlikte !!{formula}s adım adım, daha aşağıdan yukarıya da kurulabilir;
burada bunu bir
\emph{oluşum dizisi} kavramını kullanarak yapıyoruz.

\begin{defn}[Formüller için oluşum dizileri]
\ollabel{defn:fseq-frm}
$\Lang L_0$ dilinin simge dizilerinden oluşan sonlu bir
$\tuple{!A_0,\dotsc,!A_n}$ dizisi, $!A \ident !A_n$ ise ve her $i \leq n$
için ya $!A_i$ atomik bir formülse ya da aşağıdakilerden birini sağlayan
$j,k < i$ indisleri varsa $!A$ için bir 
\emph{oluşum dizisidir}:
\begin{enumerate}
    \tagitem{prvNot}{$!A_i \ident \lnot !A_j$.}{}%
    \tagitem{prvAnd}{$!A_i \ident (!A_j \land !A_k)$.}{}%
    \tagitem{prvOr}{$!A_i \ident (!A_j \lor !A_k)$.}{}%
    \tagitem{prvIf}{$!A_i \ident (!A_j \lif !A_k)$.}{}%
    \tagitem{prvIff}{$!A_i \ident (!A_j \liff !A_k)$.}{}%
\end{enumerate}
\end{defn}

\begin{ex}
\[
\tuple{
    \Obj p_0,
    \Obj p_1,
    (\Obj p_1 \land \Obj p_0),
    \lnot (\Obj p_1 \land \Obj p_0)
}
\]
$\lnot (\Obj p_1 \land \Obj p_0)$'ın bir oluşum dizisidir; şu dizi de
öyledir:
\[
\tuple{
    \Obj p_0,
    \Obj p_1,
    \Obj p_0,
    (\Obj p_1 \land \Obj p_0),
    (\Obj p_0 \lif \Obj p_1),
    \lnot (\Obj p_1 \land \Obj p_0)
}.
\]
%
İkinci örnekte görüldüğü gibi oluşum dizileri `fazlalık' içerebilir: bunlar
gereksiz olan ya da kuruluşa katkıda bulunmayan formüllerdir.
\end{ex}

\begin{prop}
\ollabel{prop:formed}
$\Frm[L_0]$'daki her !!{formula}~$!A$'nın bir oluşum dizisi vardır.
\end{prop}

\begin{proof}
$!A$'nın atomik olduğunu varsayalım. O zaman~$\tuple{!A}$ dizisi $!A$ için
bir oluşum dizisidir.
%
Şimdi $!B$ ile~$!C$'nin sırasıyla $\tuple{!B_0,\dotsc,!B_n}$ ve
$\tuple{!C_0,\dotsc,!C_m}$ oluşum dizilerine sahip olduğunu varsayalım.
%
\begin{enumerate}
  \tagitem{prvNot}{$!A \ident \lnot !B$ ise
      $\tuple{!B_0,\dotsc,!B_n,\lnot !B_n}$
      $!A$ için bir oluşum dizisidir.}{}
  \tagitem{prvAnd}{$!A \ident (!B \land !C)$ ise
      $\tuple{!B_0,\dotsc,!B_n,!C_0,\dotsc,!C_m,(!B_n \land !C_m)}$
      $!A$ için bir oluşum dizisidir.}{}
  \tagitem{prvOr}{$!A \ident (!B \lor !C)$ ise
      $\tuple{!B_0,\dotsc,!B_n,!C_0,\dotsc,!C_m,(!B_n \lor !C_m)}$
      $!A$ için bir oluşum dizisidir.}{}
  \tagitem{prvIf}{$!A \ident (!B \lif !C)$ ise
      $\tuple{!B_0,\dotsc,!B_n,!C_0,\dotsc,!C_m,(!B_n \lif !C_m)}$
      $!A$ için bir oluşum dizisidir.}{}
  \tagitem{prvIff}{$!A \ident (!B \liff !C)$ ise
      $\tuple{!B_0,\dotsc,!B_n,!C_0,\dotsc,!C_m,(!B_n \liff !C_m)}$
      $!A$ için bir oluşum dizisidir.}{}
\end{enumerate}
!!{formula}s üzerinde tümevarım ilkesi uyarınca her !!{formula} için bir
oluşum dizisi vardır.
\end{proof}

Tersini de kanıtlayabiliriz. Bu önemlidir; çünkü formülleri tanımlamaya
yönelik iki yolumuzun eşdeğer olduğunu, yani aynı sonuçları verdiğini
gösterir. Ayrıca formüller hakkındaki teoremleri, oluşum dizilerinin uzunluğu
üzerinde olağan tümevarımla kanıtlayabileceğimiz anlamına gelir.

\begin{lem}
\ollabel{lem:fseq-init}
$\tuple{!A_0,\dotsc,!A_n}$'nin $!A_n$ için bir oluşum dizisi ve $k \leq n$
olduğunu varsayalım. O zaman $\tuple{!A_0,\dotsc,!A_k}$, $!A_k$ için bir
oluşum dizisidir.
\end{lem}

\begin{proof}
Alıştırma.
\end{proof}

\begin{thm}
\ollabel{thm:fseq-frm-equiv}
$\Frm[L_0]$, $\Lang L_0$ dilinde bir oluşum dizisine sahip olan bütün simge
dizilerinin kümesidir.
\end{thm}

\begin{proof}
$F$, $\Lang L_0$ dilinde bir oluşum dizisine sahip olan bütün simge
dizilerinin kümesi olsun. \olref[pl][syn][fseq]{prop:formed}'de
$\Frm[L_0] \subseteq F$ olduğunu gördük; şimdi ters kapsamayı kanıtlayalım.

$!A$'nın $\tuple{!A_0,\dotsc,!A_n}$ oluşum dizisine sahip olduğunu
varsayalım. $!A \in \Frm[L_0]$ olduğunu son indis~$n$ üzerinde kuvvetli
tümevarımla kanıtlarız. Tümevarım varsayımımız, son indisi $r<n$ olan bir
oluşum dizisine sahip her simge dizisinin $\Frm[L_0]$'da olduğudur. Oluşum
dizisi tanımı gereği ya $!A_n$ atomiktir ya da aşağıdaki durumlardan birinin
geçerli olmasını sağlayan $j,k < n$ indisleri vardır:
\begin{enumerate}
    \tagitem{prvNot}{$!A_n \ident \lnot !A_j$.}{}%
    \tagitem{prvAnd}{$!A_n \ident (!A_j \land !A_k)$.}{}%
    \tagitem{prvOr}{$!A_n \ident (!A_j \lor !A_k)$.}{}%
    \tagitem{prvIf}{$!A_n \ident (!A_j \lif !A_k)$.}{}%
    \tagitem{prvIff}{$!A_n \ident (!A_j \liff !A_k)$.}{}%
\end{enumerate}
Şimdi durumlara göre akıl yürütelim. $!A_n$ atomikse
$!A_n \in \Frm[L_0]$'dır. Bunun yerine $!A \ident
(!A_j \land !A_k)$ olduğunu varsayalım.
\olref[pl][syn][fseq]{lem:fseq-init} uyarınca
$\tuple{!A_0,\dotsc,!A_j}$ ve $\tuple{!A_0,\dotsc,!A_k}$ sırasıyla $!A_j$
ve~$!A_k$ için oluşum dizileridir. Bunların son indisleri $j,k<n$'dir
(dolayısıyla uzunlukları $n+1$'den küçüktür). Bu nedenle tümevarım varsayımı
uyarınca $!A_j$ ve~$!A_k$, $\Frm[L_0]$'dadır; öyleyse !!a{formula} olarak tanımlanma kuralı
gereği $(!A_j \land !A_k)$ da $\Frm[L_0]$'dadır. Öteki durumlar paralel bir
akıl yürütmeyle çıkar.
\end{proof}

\end{document}
